\documentclass[10pt,a4paper]{amsart}
\usepackage[margin=19mm]{geometry}
\usepackage{amsmath,amssymb,mathtools}
\usepackage[scr=rsfs,cal=euler]{mathalfa}
\usepackage{xfrac}
\usepackage[unicode,hidelinks,bookmarks=false]{hyperref}
\newcommand{\ie}{i.e.\ }
\newcommand{\cf}{cf.\ }
\newcommand{\resp}{respectively\ }
\newcommand{\Subsection}{\subsection}
\providecommand{\nobreakdash}{\nobreak\hbox{-}}
\setlength{\emergencystretch}{2em}
\multlinegap0pt
\usepackage{bm}
\usepackage{bbm}
\usepackage{dsfont}
\usepackage{xspace}
\usepackage{cancel}
\usepackage{mathtools}
\usepackage{amsthm}
\makeatletter
\@ifundefined{defi}{\newtheorem{defi}{Defi}}{}
\@ifundefined{thm}{\newtheorem{thm}{Thm}}{}
\makeatother
\newcommand{\Act}{\mathsf{Act}}
\newcommand{\AtAct}{\mathsf{At}}
\newcommand{\CoOp}{\mathsf{Op}}
\newcommand{\Formu}{\mathsf{Form}}
\newcommand{\LTDyn}{\mathcal{L}_{\func{F}, \Pred}^\mathsf{Dyn}}
\newcommand{\LT}{\mathcal{L}_{\func{F},\Pred}}
\newcommand{\NMH}{\mathcal{H}} 
\newcommand{\Pred}{\Lambda}
\newcommand{\TeOp}{\mathsf{Te}}
\newcommand{\alg}[1]{\mathbf{#1}}
\newcommand{\arity}{\mathsf{ar}}
\newcommand{\diam}[1]{\langle {#1} \rangle}
\newcommand{\ev}{\mathsf{ev}}
\newcommand{\func}[1]{\mathsf{#1}}
\newcommand{\nkTempl}{\mathsf{Templ}_{n,k}}
\newcommand{\sem}[1]{[\![{#1}]\!]}
\newcommand{\test}{\mathsf{test}}
\title[Many-valued coalgebraic dynamic logics: Safety and strong co]{Many-valued coalgebraic dynamic logics: Safety and strong completeness via reducibility}
\author{Helle Hvid Hansen, Wolfgang Poiger}
\date{Source: arXiv:2512.22851v2 (CC BY 4.0)}
\begin{document}
\maketitle

\noindent\emph{Source and licence.} Many-valued coalgebraic dynamic logics: Safety and strong completeness via reducibility. Version arXiv:2512.22851v2 (CC BY 4.0), distributed under the Creative Commons Attribution 4.0 International licence, as recorded in the arXiv licence metadata for this exact version and shown on the abstract page https://arxiv.org/abs/2512.22851v2. Full rights notice: licenses/arxiv-2512.22851-CC-BY-4.0.txt.\par\smallskip

\noindent\textit{Editorial scope.} Counted unit: the Thm whose source label is \texttt{theorem:StandardQuasiCanModel} in the article (source0.tex lines 1487--1493, source label \texttt{theorem:StandardQuasiCanModel}), delivered together with the article's own complete original proof of it (source0.tex lines 1495--1545, one \texttt{proof} environment of 51 lines). No mathematics is written, repaired or re-derived: statement and proof are the article's own text line for line. Required context retained uncounted: definition:StandardModel (lines 688--690); definition:ReducibleCoalgOperation (lines 994--1000); definition:Reducibletest (lines 1303--1308); definition: LTDyn (lines 1435--1446). Every definition, display and label of those lines is retained unchanged. Omitted from this package and not redistributed: 1--687, 691--993, 1001--1302, 1309--1434, 1447--1486, 1494--1494, 1546--1811. Nothing inside a retained excerpt is omitted or altered: every retained body is delivered exactly as the article's lines are written, so no omission receipt of any kind accompanies this package. Citations: the retained excerpts carry no citation command at all, so no bibliography entry of the article is transcribed and no reference is redistributed. Reference adaptation: none was needed and none was made; every reference inside the delivered unit resolves inside it, so no unresolved reference or citation is produced. Printed slips kept literal: no printed word, symbol, hypothesis or number of the counted statement or of its proof was altered. Layout and class: the article's own class is \texttt{lmcs} with packages including amsmath, amsthm, amssymb, tikz, graphicx, xy, tikz-cd, hyperref, leftindex, comment, euflag; the wrapper is the lane's pinned \texttt{amsart} editorial wrapper at 10pt on A4, which restates the theorem-like environments the retained text needs and adds the article's own macro definitions that the retained text needs (\texttt{\textbackslash Act}, \texttt{\textbackslash AtAct}, \texttt{\textbackslash CoOp}, \texttt{\textbackslash Formu}, \texttt{\textbackslash LT}, \texttt{\textbackslash LTDyn}, \texttt{\textbackslash NMH}, \texttt{\textbackslash Pred}, \texttt{\textbackslash TeOp}, \texttt{\textbackslash alg}, \texttt{\textbackslash arity}, \texttt{\textbackslash diam}, \texttt{\textbackslash ev}, \texttt{\textbackslash func}, \texttt{\textbackslash nkTempl}, \texttt{\textbackslash sem}, \texttt{\textbackslash test}), with extra wrapper packages (bm, bbm, dsfont, xspace, cancel, mathtools). The delivered numbering is the unit's own. Assets: no retained span contains a figure environment, a graphics-inclusion command, a PDF-page inclusion, a table or a TikZ picture; no image file is copied into this package, the package declares no asset dependency and no image file is redistributed with it. Licence: version arXiv:2512.22851v2 of this article is distributed under the Creative Commons Attribution 4.0 International licence, as recorded in the arXiv licence metadata for this exact version and shown on the abstract page https://arxiv.org/abs/2512.22851v2. A full rights notice is delivered at licenses/arxiv-2512.22851-CC-BY-4.0.txt. Characters: every retained excerpt is pure ASCII, so the delivered engine, XeLaTeX with its default Latin Modern text font, reports no missing character.
\begin{defi}[\textsf{Standard models}]\label{definition:StandardModel}
A model $\gamma \colon \Act \to (\func{F}X)^X$ is \emph{standard} if it is obtained from some $\gamma^0 \colon \AtAct \to (\func{F}X)^X$ together with a propositional valuation as described above. 
\end{defi} 

\begin{defi}[\textsf{Reducible coalgebra operations}]\label{definition:ReducibleCoalgOperation}
Let $O\in \CoOp$ be an $n$-ary coalgebra operation and $\lambda \in \Pred$ be a $k$-ary predicate lifting. We call $O$ \emph{reducible with respect to $\lambda$} if there exists an $(n, k)$-template $\tau \in \nkTempl$ such that 
$$
\lambda_X(\vec{\sigma}) \circ O(\vec{\gamma}) = \tau[\vec{\gamma}, \vec{\sigma}]
$$   
holds for all $\vec{\gamma}\colon X {\to} \func{F}^n X$ and $\vec{\sigma} \in  (A^X)^k $. We simply call $O$ \emph{reducible} if it is reducible with respect to \emph{every} $\lambda \in \Pred$. 
\end{defi} 

\begin{defi}[\textsf{Reducible tests}]\label{definition:Reducibletest}
Let $\test \in \TeOp$ be a test and $\lambda \in \Pred$ be a $k$-ary predicate lifting. We call $\test$ \emph{reducible with respect to $\lambda$} if there is a $(k+1)$-ary $\alg{A}$-term function $t(x, x_1, \dots, x_k) \colon A^{k+1} \to A$ such that 
$$
\lambda_X(\sigma_1, \dots, \sigma_k) \circ \test(\sigma) = t^\mathrm{pw}(\sigma, \sigma_1, \dots, \sigma_k)$$
for all $\sigma, \sigma_i \in A^X$. We call $\test$ \emph{reducible} if it is reducible with respect to \emph{every} $\lambda \in \Pred$.  
\end{defi} 

\begin{defi}[\textsf{Coalgebraic dynamic logic $\LTDyn$}]\label{definition: LTDyn}
The \emph{dynamic coalgebraic logic} $\LTDyn \subseteq \Formu$ is the smallest set of dynamic formulas satisfying the following conditions. 
\begin{itemize}
    \item For every action $a \in \Act$, the set $\LTDyn$ contains the axioms of $\LT$ with respect to the modalities $\{\diam{a}^\lambda\}_{\lambda\in \Lambda}$ and is closed under the corresponding rules of $\LT$.
    \item For every pair $(O, \lambda) \in \CoOp\times\Pred$, the set $\LTDyn$ contains all instances of the reduction axiom 
 $$\diam{O(\vec{a})}^\lambda(\vec{\varphi}) \leftrightarrow \tau[\vec{a}, \vec{\varphi}],$$
where the template $\tau$ witnesses reducibility of $O$ with respect to $\lambda$ (Definition~\ref{definition:ReducibleCoalgOperation}).
    \item For every pair $(\test, \lambda) \in \TeOp\times\Pred$, the set $\LTDyn$ contains all instances of the reduction axiom 
    $$\diam{\test(\psi)}^\lambda (\vec{\varphi}) \leftrightarrow t(\psi, \vec{\varphi}),$$
where the $\alg{A}$-term-function $t$ witnesses reducibility of $\test$ with respect to $\lambda$ (Definition~\ref{definition:Reducibletest}).
\end{itemize}   
\end{defi}

\begin{thm}\label{theorem:StandardQuasiCanModel}
Let $\zeta^0 \colon \AtAct \to (\func{F}\NMH)^{\NMH}$ be coherent for \emph{atomic} actions, that is, 
$$
\lambda_\NMH \big(\ev_{\varphi_1}, \dots, \ev_{\varphi_{\arity(\lambda)}}\big)\big(\zeta_{a_0}^0 (h)\big) = h\big(\diam{a_0}^\lambda (\varphi_1, \dots, \varphi_{\arity(\lambda)})\big) 
$$
holds for all $a_0 \in \AtAct$, $\lambda \in \Pred$, $\varphi_i \in \Formu$ and $h \in \NMH$. Let $\zeta \colon \Act \to (\func{F}\NMH)^\NMH$ be the standard extension of $\zeta^0$ as in Definition~\ref{definition:StandardModel}. Then $\zeta$ is a quasi-canonical standard model for $\LTDyn$.     
\end{thm} 

\begin{proof}
By mutual induction we show that $\sem{\varphi}^c = \ev_{\varphi} \colon \NMH \to A$ for all formulas $\varphi \in \Formu$ and coherence (as in the statement) for all actions $a \in \Act$ hold in $\zeta$. Here, only the case $\varphi = \diam{a}^\lambda (\varphi_1, \dots, \varphi_k)$ for some non-atomic action $a \in \Act$ and $k$-ary $\lambda \in \Pred$ is not immediate.  

\textbf{Coalgebra operations.} Let $a = O(a_1, \dots ,a_n)$ for some $n$-ary coalgebra operation $O \in \CoOp$ and inductively assume that coherence already holds for all $\zeta_{a_j}$. Let $\tau$ be the $(n,k)$-template corresponding to the reduction axiom for $(O, \lambda)$ which is contained in $\LTDyn$ (recall Definition~\ref{definition: LTDyn}). Since this reduction axiom is contained in $\LTDyn$, for every $h \in \NMH$ we have  
$$
h\big(\diam{O(a_1, \dots, a_n)}^\lambda (\varphi_1, \dots, \varphi_k)\big) = h\big(\tau[a_1, \dots, a_n, \varphi_1, \dots, \varphi_k]\big),$$ 
which shows $\ev_{\diam{O(a_1, \dots, a_n)}^\lambda (\varphi_1, \dots, \varphi_k)} = \ev_{\tau[a_1, \dots, a_n, \varphi_1, \dots, \varphi_k]}$. On the other hand, since $\tau$ witnesses reducibility, together with the inductive hypothesis on $\varphi_1, \dots, \varphi_k$ we get 
\begin{align*}
\sem{\diam{O(a_1, \dots, a_n)}^\lambda (\varphi_1, \dots, \varphi_k)}^\mathrm{c} & = \lambda_\NMH (\ev_{\varphi_1}, \dots, \ev_{\varphi_k}) \circ O(\zeta_{a_1}, \dots, \zeta_{a_n}) \\ 
& = \tau[\zeta_{a_1}, \dots, \zeta_{a_n}, \ev_{\varphi_1}, \dots, \ev_{\varphi_k}].
\end{align*}
Thus, it suffices to show that for every $\tau \in \nkTempl$ it holds that 
$$\ev_{\tau[\vec{a}, \vec{\varphi}]} = \tau[\zeta_{\vec{a}}, \ev_{\vec{\varphi}}],$$  
where we abbreviate $\zeta_{\vec{a}} = (\zeta_{a_1}, \dots, \zeta_{a_n})$ and $\ev_{\vec{\varphi}} = (\ev_{\varphi_1}, \dots, \ev_{\varphi_k})$. To show this, we proceed by induction on the structure of the reduction template $\tau$. 

For $\tau = \omega_i$, directly by definition we get 
$$
\ev_{\omega_i[\vec{a}, \vec{\varphi}]} = \ev_{\varphi_i} = \omega_i[\zeta_{\vec{a}}, \ev_{\vec{\varphi}}]
$$
as desired.

For $\tau = o(\tau_1 , \dots, \tau_{\arity(o)})$ we use the inductive hypothesis on $\tau_i$ for all $i = 1, \dots, \arity(o)$ to compute
$$
\tau[\zeta_{\vec{a}}, \ev_{\vec{\varphi}}] = o^\mathrm{pw} \big(\tau_1[\zeta_{\vec{a}}, \ev_{\vec{\varphi}}], \dots, \tau_{\arity(o)}[\zeta_{\vec{a}}, \ev_{\vec{\varphi}}]\big) = o^\mathrm{pw}(\ev_{\tau_1[\vec{a}, \vec{\varphi}]}, \dots, \ev_{\tau_{\arity(o)}[\vec{a}, \vec{\varphi}]}).
$$
Using the fact that $h$ (being a non-modal homomorphism) commutes with the $\alg{A}$-operation $o$ we observe that the above sends every $h \in \NMH$ to  
$$
o^\mathrm{pw}(\ev_{\tau_1[\vec{a}, \vec{\varphi}]}, \dots, \ev_{\tau_{\arity(o)}[\vec{a}, \vec{\varphi}]})(h) = h\big (o(\tau_1[\vec{a}, \vec{\varphi}], \dots, \tau_{\arity(o)}[\vec{a}, \vec{\varphi}]) \big) = h\big(\tau[\vec{a}, \vec{\varphi}]\big)
$$
as desired. 

Lastly, for the case $\tau = \diam{j}^{\lambda'}(\tau_1, \dots, \tau_{\arity(\lambda')})$, first we use the inductive hypothesis on the $\tau_i$ for all $i = 1,\dots, \arity(\lambda)$ and get 
\begin{align*}
\tau[\zeta_{\vec{a}}, \ev_{\vec{\varphi}}] & = \lambda_\NMH'\big(\tau_1[\zeta_{\vec{a}}, \ev_{\vec{\varphi}}], \dots, \tau_{\arity(\lambda')}[\zeta_{\vec{a}}, \ev_{\vec{\varphi}}]\big) \circ \zeta_{a_j} \\
& = \lambda_\NMH'(\ev_{\tau_1[\vec{a}, \vec{\varphi}]}, \dots, \ev_{\tau_{\arity(\lambda')}[\vec{a}, \vec{\varphi}]} ) \circ \zeta_{a_j}.
\end{align*}
Now we make use of the inductive hypothesis on coherence with respect to the action $a_j$, which yields that for every $h \in \NMH$ we have 
$$
\lambda_\NMH'(\ev_{\tau_1[\vec{a}, \vec{\varphi}]}, \dots, \ev_{\tau_{\arity(\lambda')}[\vec{a}, \vec{\varphi}]} ) \big( \zeta_{a_j}(h)\big) = h\big(\diam{a_j}^{\lambda'} (\tau_1[\vec{a}, \vec{\varphi}], \dots, \tau_{\arity(\lambda')}[\vec{a}, \vec{\varphi}])\big)
$$
as desired. This finishes the case of composite actions obtained via coalgebra operations.

\textbf{Tests.}
The last case we need to discuss is $a = \test(\psi)$, inductively assuming that $\sem{\psi} = \ev_\psi$ already holds. Let $t(x, x_1, \dots, x_k) \colon A^{k+1} \to A$ be the $\alg{A}$-term-function corresponding to the reduction axiom for $(\test, \lambda)$ which is included in $\LTDyn$. Then for any $h \in \NMH$ we get  
\begin{align*}
\sem{\diam{\test(\psi)}^\lambda(\vec{\varphi})}^\mathrm{c} & = \lambda_{\NMH}(\ev_{\varphi_1}, \dots, \ev_{\varphi_k})\big(\test(\ev_{\psi})(h)\big) && \text{(Def. \& Ind. Hyp.)} \\ 
& = t\big(\ev_{\psi}(h), \ev_{\varphi_1}(h), \dots, \ev_{\varphi_k}(h)\big) && \text{(Reducibility $\test$)} \\ 
& = t\big(h(\psi), h(\varphi_1), \dots,h(\varphi_k)\big)  = h\big(t(\psi, \vec{\varphi})\big). && \text{(Definitions $\ev$ \& $h \in \NMH$)}  
\end{align*} 
Finally, we have $h(t(\psi, \vec{\varphi})) = h(\diam{\test(\psi)}\vec{\varphi})$ since this reduction axiom is included in $\LTDyn$.   
\end{proof}

\end{document}
